%%%%%%%%%%%%%%%%%%%%%%%%%%%%%%%%%%%%%%%%%%%%%%%%%%%%%%%%%%%%%%%%%%%%%%%%%%%%%%%
\section{Introduction}\label{sec:intro}
%%%%%%%%%%%%%%%%%%%%%%%%%%%%%%%%%%%%%%%%%%%%%%%%%%%%%%%%%%%%%%%%%%%%%%%%%%%%%%%
The use of modern computer systems almost always requires that a user prove their identity through some process of authentication.  Specifically, a user can authenticate using methods such as public key authentication, biometrics, and passwords; something you have, are, or know, respectively.  Password-based authentication remains the most widely used option for authentication because of its ease of use and simple design. %TODO: Citation here?

%this is the original of the paragraph below
%However, passwords have a human factor weakness.  Users often choose passwords that are too simplistic and easily guessed.  Thus, administrators encourage or force users to select more complex and longer passwords.  More complex and longer passwords decrease user satisfaction because entering the passwords is perceived as difficult.  This difficulty is compounded when the user enters a password incorrectly and has to re-enter it.  Additionally, the likelihood of entering a password incorrectly only increases with password complexity and length as does the time to enter and re-enter it.  The difficulty of entering a complex password rises to an unacceptable level very quickly in a medical or clinical setting, particularly when patient needs are urgent.  

However, passwords have a human factor weakness as users often choose passwords that are too simplistic and easily guessed~\cite{Yan:2004:PMS:1024867.1025014}.  Administrators and systems require users to select more complex passwords as a consequence; thus, decreasing user satisfaction as password selection becomes seemingly difficult. Worst yet, complex passwords interfere in critical workflow such as clinical care where a patient's need is most urgent.  

In this paper we describe an authentication system that requires the user to enter a password as infrequently as once a day.  Specifically, authentication information is stored on a wearable device, a bracelet in our case, and is transmitted to devices to which the user wishes to authenticate. The transmission between the bracelet and device is achieved via using the user's body as a communication medium~\cite{Barth:2008:BCB:1460257.1460273,chang:healthtech:2013}. Our goal is to reduce the impact on user satisfaction and workflow by removing most of the difficulty of using a complex password. 

While we focus on authentication in the medical community use case, we anticipate that other areas such as the financial sector may benefit from our system. In addition, it is important to note that this system is not a two-factor authentication solution. The bracelet is not a biometric component and does not provide any additional information outside of what it stores. It is meant to enhance the user experience and encourage the use of complex passwords. Finally, KBID is a work in progress.

%Our goal is to authenticate a user in less than one second provided that user has provided valid credentials once.   The hope is to reduce the resistance to strong password requirements by removing the vast majority events where the user is asked to enter a password.  %To put it simply, users will be more willing to use a strong password if they only have to enter it once a day as opposed to being prompted to enter the password each time they use a computer system.

%While we have focused on developing an authentication system for the medical community, there is nothing in the authentication system that prevents it from being used in other areas.  We anticipate that the financial sector may find our system interesting.

%It is important to note what this system is not.  This system is not a two factor authentication solution.  The bracelet should not be confused as the provider of additional authentication information outside of what it is instructed to hold.  There is no biometric component to this system.  This system is not a complete authentication system.  It is meant to be an enhancement to the user experience and to encourage the use of strong passwords by eliminating the need to repeatedly entering them. Finally, KBID is a work in progress.



%To use the system, the user will don a wearable device.  Currently we envision this device as a bracelet.  The user will approach a host computer equipped with a special hardware to communicate with the bracelet.  We will refer to this hardware as an authentication module.  The user will touch a sensor on the authentication module and the authentication module will communicate over the user's skin and determine that the user is not authenticated.  The authentication module will then communicate with a software client on the host computer and the software will prompt the user to log in.  The user will provide a username and password that will be validated with the back end authentication server.  If the user's credentials are valid, the server will provide authentication information to the client.  The client will instruct the user to touch the sensor on the authentication module again.  When the user touches the authentication module the client will send the authentication information to the authentication module which will in turn pass it over the user's skin to the bracelet where it will be stored.

%The next time the user wishes to authenticate to a host computer equipped with this system, the user will touch the authentication module.  The authentication module will ask the bracelet for the authentication information and the bracelet will provide it.  The authentication module will then send the authentication information to the client on the host computer.  The client will verify the validity of the authentication information with the authentication server.  If the authentication information is valid the client will unlock the host computer.  It is these subsequent authentication events that should take less than one second.

%If the user removes the bracelet, the bracelet has a means of detecting its removal.  When it detects its removal the bracelet will delete the authentication information it contains.  In order to continue using the system, the user must put the bracelet back on and re-authenticate using a username and a password.

%It is conceivable that users of computer systems in a clinical setting may not be able to use a keyboard and mouse due to sanitation concerns.  An additional benefit of this system is that clinicians could authenticate by making contact between the sensor and another part of their arm, such as the elbow, without compromising the sanitization of their hands.

%While we have focused on developing an authentication system for the medical community, there is nothing in the authentication system that prevents it from being used in other areas.  We anticipate that the financial sector may find our system interesting.

%At their developer conference in August 2015, Intel announced a wearable authentication solution.  The Intel solution has a very similar workflow and seems to achieve similar results.  However the Intel solution uses Bluetooth Low Energy (BLE).  We speculate that the Intel solution may be vulnerable to a range extending attack.  In the speculative attack, an attacker intercepts the BLE signal from the user's wearable device and the computer to which they are authenticated.  The attacker then rebroadcasts the BLE traffic over a longer distance, possibly by translating the traffic to another section of the RF spectrum during the retransmission process.  No authentication or paring is required since all the attacker is doing is receiving and then rebroadcasting BLE traffic from previously paired and authenticated equipment.  Our system reduces or eliminates this type of attack by requiring that the user be in physical contact with the authentication device to transmit authentication information.

%this is the original of the paragraph below
%It is important to note what this system is not.  This system is not a two factor authentication solution.  It is conceivable that this system could include a two factor authentication component but the bracelet should not be confused as the provider of additional authentication information outside of what it is instructed to hold.  There is no biometric component to this system.  Neither the bracelet nor the authentication module measure or store any information about the user's physical state.  In this system the user's body is simply the medium across which the authentication information is communicated.  This system is not a complete authentication system.  It is meant to be an enhancement to the user experience and to encourage the use of strong passwords by eliminating the need to repeatedly entering them.  As such there will be no discussion in the paper about access control, password hashing, and much of the intricacies of the of the authentication process.  For this research Kerberos was used as the authentication provider but it is important to note that this system will support the use of any authentication provider that requires the use of stored tokens or tickets.  It could also be enhanced to provide its own tokens or tickets, though that is beyond the scope of the planned research on this project.  Finally, KBID is a work in progress.
